\section{Arbeit im Regelungstechniklabor}

\subsection{Rolle der Betreuer\_innen}
Deine Erst- und Zweitprüfer\_innen sind erste Ansprechpartner\_innen bei allen auftretenden Unklarheiten. Neben fachlichen Fragen werden sie dir auch bei der Beantwortung aller anderen Fragen, die im Zusammenhang mit deiner Abschlussarbeit an der HTW auftreten, z. B. Beschaffungen, Arbeitsaufträge an die Werkstätten, etc., helfen,  Die Initiative sollte jedoch stets von dir ausgehen.


\subsection{Arbeitsplatz, Arbeitszeit und Schlüssel im Regelungtechnik- Labor}

Wenn benötigt, lass dir zu Beginn deiner Arbeit vom Laboringenieur einen Arbeitsplatz und, wenn vorhanden, ein abschließbares Schrankfach zuweisen und die Schlüsselregeln  erklären. Bitte akzepiere zu Beginn deiner Arbeit im Labor die aktuellen Hygieneregeln und die Laborordnung mit deiner Unterschrift.

Die Arbeitstische im Regelungstechnik- Labor werden in der Regel von mehreren Personen benutzt. Räumen Sie daher Ihren Tisch nach Benutzung stets auf.

Die Arbeitsräume sind während der regulären Arbeitszeit von 8:00 bis 20:00 Uhr zugänglich. Die Laborordung ist Um dir die Möglichkeit zu geben, in begründeten Ausnahmefällen auch in der Zeit von 20.00 bis 22.00 Uhr noch arbeiten zu können, wird folgende Regelung vereinbart:
\begin{itemize}
    \item In diesem Zeitraum dürfen ohne Betreuer\_in keine Hardware-Arbeiten durchgeführt werden.
    \item Das Haus muss spätestens kurz vor 22.00 Uhr verlassen werden. Ab 22.00 Uhr werden die Haustüre und die Flureingangstüren vom Schließdienst versperrt.
    \item Jeder außerhalb der regulären Arbeitszeit und ohne Betreuer arbeitende Diplomand ist verantwortlich für alle selbst verursachten Schäden. Alle Haftungsansprüche gegenüber dem  Lehrstuhl und den für den Lehrstuhl handelnden Personen sind, soweit gesetzlich zugelassen, ausgeschlossen.
\end{itemize}


\subsection{Werkzeug, Werkstattmaterial, Werkstattaufträge}

Werkzeuge und Werkstattmaterial werden vom Laboringenieur ausgegeben. Werkstattaufträge sind mit entsprechenden Konstruktionszeichnungen/-skizzen mit der Betreuerin/ dem Betreuer abzuklären.

\subsection{Laborgeräte, Gerätebeschreibungen und Bauelemente}

Laborgeräte und Gerätebeschreibungen können Sie beim LÖaboringenieur entleihen. 
%Im Gerätelager muss ein Entnahmeschein  hinterlegt werden, der vom Betreuer zu unterzeichnen ist. 
Alle nicht mehr benutzten Geräte sind sofort in das Lager zurückzugeben. Achten Sie bitte auf eine sachgemäße und pflegliche Behandlung der Geräte. Alle Schäden an den Geräten sind umgehend schriftlich an den Laboringenieur zu melden. Wenn fahrlässiges Verschulden vorliegt, wird die Benutzerin/ der Benutzer zum Schadenersatz herangezogen.

Elektronische Bauelemente werden auch vom Laboringenieur ausgegeben.

\subsection{Einkäufe und Bestellungen}

Einkäufe und Bestellungen sind über den Laboringenieur abzuwickeln. Werden Sie mit Besorgungen betraut, so achten Sie bitte sorgfältig auf Kassenbelege, Lieferscheine und Rechnungen, die alle sofort bei Ihrem Betreuer abzuliefern sind. 

\subsection{Haftung}
Die Haftung des/der Diplomanden/Diplomandin erstreckt sich auf alle Leihgaben, die im abschließbaren Schrankfach untergebracht werden können. 
 %Wir empfehlen  dringend, Werkzeuge, kleinere Messgeräte, Bücher und Bauteile nach Beendigung der Arbeit stets wegzuschließen.

